\documentclass[conference,a4paper,10pt]{IEEEtran}
\IEEEoverridecommandlockouts
\usepackage[utf8]{inputenc}
\usepackage{lmodern}
\usepackage{amsmath,amssymb}
\usepackage{graphicx}
\usepackage{booktabs}
\usepackage{array}
\usepackage{xcolor}
\usepackage{hyperref}
\usepackage{listings}
\usepackage{caption}
\usepackage{float}
\usepackage{multirow}

\definecolor{codegray}{RGB}{245,245,245}
\definecolor{codegreen}{RGB}{0,110,0}
\definecolor{codeblue}{RGB}{0,0,180}
\definecolor{codered}{RGB}{160,0,0}

\lstset{
  backgroundcolor=\color{codegray},
  basicstyle=\ttfamily\scriptsize,
  keywordstyle=\color{codeblue}\bfseries,
  commentstyle=\color{codegreen}\itshape,
  stringstyle=\color{codered},
  breaklines=true, frame=single,
  numbers=left, numberstyle=\tiny\color{gray},
  tabsize=4, language=Python
}

\hypersetup{colorlinks=true,linkcolor=black,citecolor=black,urlcolor=black}

\begin{document}

% ─────────────────────────────────────────────────────────────
\title{\textbf{Workshop No.~4 --- System Simulation and Validation}\\
{\large Community Security Alert System (CSAS)}}

\author{
\IEEEauthorblockN{%
  Felipe Jose Garzon Herrera$^{1}$,
  Juan Esteban Quintero Gordillo$^{2}$,
  Henry Samuel Garrido Medina$^{3}$,
  Gabriel Mateo Cusba Marin$^{4}$}\\
\IEEEauthorblockA{Systems Analysis \& Design --- Semester 2026-I\\
Computer Engineering Program,
Universidad Distrital Francisco Jos\'e de Caldas\\
Instructor: Eng.\ Carlos Andr\'es Sierra, M.Sc.\\
$^{1}$Cod.~20251020132\quad
$^{2}$Cod.~20251020137\quad
$^{3}$Cod.~20251020125\quad
$^{4}$Cod.~20251020128}
}
\maketitle

% ─────────────────────────────────────────────────────────────
\begin{abstract}
This report presents the simulation and validation of the Community
Security Alert System (CSAS), a platform designed to improve campus
security at Universidad Distrital Francisco Jos\'e de Caldas. Building
on the analysis from Workshops 1, 2, and 3, this workshop implements
two complementary simulation approaches: a Discrete Event Simulation
(DES) using the SimPy library, and a hybrid behavioral simulation that
captures community adoption dynamics and feedback loops. Three
operational scenarios are evaluated---Baseline, Optimization, and
Failure Mode---each executed 30 independent times for statistical
reliability. The results show that the Optimization scenario meets
the main performance targets: average latency of 17.94~s (well below
the 30~s threshold), an alert delivery success rate of 96.84\,\%, and
a community adoption rate of 67.62\,\%, exceeding the 60\,\% goal.
The most important finding is that the verification threshold is
the variable that most influences system stability.
\end{abstract}

% ─────────────────────────────────────────────────────────────
\section{Introduction}

The purpose of this workshop is to evaluate the operational behavior of
the CSAS under different campus security conditions through
computational simulation.

Throughout the previous workshops, the team identified the main
security problems at the university (Workshop~1), designed the system
architecture to address them (Workshop~2), and defined the risk
management and project plan (Workshop~3). This workshop takes that
design and tests it in a controlled simulation environment to verify
whether the proposed architecture actually works as expected.

The simulation models the complete incident reporting workflow,
including incident generation, AI-based verification, geofenced alert
broadcasting, emergency prioritization, and performance monitoring.
Two simulation approaches were used because each one captures a
different aspect of the system: the DES focuses on the technical
event pipeline, while the behavioral simulation focuses on how the
community responds over time.

% ─────────────────────────────────────────────────────────────
\section{Simulation Paradigm}

\subsection{Discrete Event Simulation (DES)}

The DES approach was chosen for the process-oriented part of the
simulation because the CSAS platform operates through asynchronous
events that happen at specific moments in time. Instead of modeling
the system continuously, the simulation advances from event to event,
which makes it efficient for modeling queues, processing delays,
concurrent users, and resource utilization~\cite{law2015simulation}.

Examples of events modeled in the DES are: a user submitting an
incident report, the verification service processing that report,
the dispatcher broadcasting an alert, and the analytics module
recording the response time.

The DES was implemented in Python using the \texttt{SimPy}
library~\cite{simpy}, which provides the core event queue and
process management.

\subsection{Behavioral Simulation}

The behavioral simulation complements the DES by modeling something
the event queue cannot capture: how the community behaves over time.
Two feedback loops govern long-term system behavior:

\begin{itemize}
  \item \textbf{Reinforcing loop (R+):} when the system correctly
    verifies and broadcasts an alert, community trust and vigilance
    increase, which leads to more reports in future steps---a
    self-sustaining cycle.
  \item \textbf{Balancing loop (B-):} when too many unfiltered or
    false alerts are broadcast, users experience alert fatigue. Once
    fatigue exceeds 60\,\%, adoption drops by 4 percentage points
    per step, eventually collapsing participation.
\end{itemize}

Which loop dominates depends almost entirely on the verification
threshold parameter, which is the core finding of the behavioral
analysis.

% ─────────────────────────────────────────────────────────────
\section{System Architecture}

The simulated architecture reflects the design developed in
Workshop~2, organized around five independent modules. Each module
corresponds to a microservice in the real system.

\begin{table}[H]
\caption{Simulated System Modules}
\label{tab:modules}\centering\small
\begin{tabular}{ll}
\toprule
\textbf{Module} & \textbf{Responsibility}\\
\midrule
User Process        & Generates incidents dynamically\\
Incident Service    & Receives and stores reports in a queue\\
Verification Service & Simulates AI and community verification\\
Dispatcher Service  & Sends geofenced alerts to nearby users\\
Analytics Module    & Collects performance metrics\\
\bottomrule
\end{tabular}
\end{table}

The verification logic applies the following condition before
broadcasting any alert:
\[
\text{verified} = (\text{confidence} \geq 0.75)
\;\land\;
(\text{confirmations} \geq \theta \;\lor\; \text{admin\_approval})
\]
where $\theta$ is the verification threshold and \textit{confidence}
is a randomly sampled AI score. This condition directly models
FR-03 (AI Verification) and NFR-08 (Data Integrity).

% ─────────────────────────────────────────────────────────────
\section{Functional Requirements Representation}

Table~\ref{tab:fr} shows how each functional requirement from
Workshop~2 was represented in the simulation.

\begin{table}[H]
\caption{Functional Requirements Coverage}
\label{tab:fr}\centering\small
\begin{tabular}{ll}
\toprule
\textbf{FR} & \textbf{Simulation Representation}\\
\midrule
FR-01 & Users generate security incidents\\
FR-02 & Alerts sent within a 500\,m geofence radius\\
FR-03 & Confidence-based AI validation ($\geq 0.75$)\\
FR-04 & Simulated through console monitoring and CSV export\\
FR-05 & Community confirmations randomly generated [0,\,4]\\
FR-06 & Random GPS coordinates assigned per incident\\
FR-07 & Incidents stored in a SimPy in-memory queue\\
FR-08 & SOS flag enables emergency priority processing\\
FR-09 & Metrics and performance indicators generated per step\\
FR-10 & Reports may include a simulated media attachment flag\\
\bottomrule
\end{tabular}
\end{table}

% ─────────────────────────────────────────────────────────────
\section{Non-Functional Requirements Representation}

\begin{table}[H]
\caption{Non-Functional Requirements Coverage}
\label{tab:nfr}\centering\small
\begin{tabular}{ll}
\toprule
\textbf{NFR} & \textbf{Simulation Mechanism}\\
\midrule
NFR-01 & Response time measured per incident; target $< 30$~s\\
NFR-02 & Dispatcher failure probability $p = 2\,\%$\\
NFR-03 & Continuous execution over the full simulation window\\
NFR-04 & Conceptually represented (TLS / AES at infrastructure level)\\
NFR-05 & Multiple concurrent users simulated\\
NFR-06 & Simplified $\leq 3$-interaction reporting process\\
NFR-07 & Modular class architecture (one class per service)\\
NFR-08 & Dual-condition verification gate enforced\\
NFR-09 & Multi-confirmation validation threshold $\theta$\\
\bottomrule
\end{tabular}
\end{table}

% ─────────────────────────────────────────────────────────────
\section{Simulation Logic}

The DES simulation follows a sequential workflow for each incident:

\begin{enumerate}
  \item A user process generates an incident at a random interval.
  \item The incident is placed in the Incident Service queue.
  \item The Verification Service evaluates it using the AI confidence
    score and community confirmation count.
  \item If the verification condition is met, the incident is marked
    as verified.
  \item The Dispatcher Service broadcasts the alert to users within
    the 500\,m geofence.
  \item The response time is recorded from incident creation to
    successful dispatch.
  \item All metrics are collected for the final analysis.
\end{enumerate}

In the behavioral simulation, these steps are abstracted into a
per-step update of system state variables: adoption rate, fatigue
level, trust level, and vigilance. Each step represents approximately
one hour of real operation (window: 15:00--02:00), with environmental
multipliers applied based on time of day, weather conditions, and
campus zone.

% ─────────────────────────────────────────────────────────────
\section{Simulation Parameters}

\begin{table}[H]
\caption{Simulation Configuration Parameters}
\label{tab:params}\centering\small
\begin{tabular}{ll}
\toprule
\textbf{Parameter} & \textbf{Value}\\
\midrule
Simulation time (DES)       & 200 seconds\\
Simulation steps (behav.)   & 12 steps (15:00--02:00 window)\\
Number of users (DES)       & 11\\
Population (behavioral)     & 200 users (pilot scope)\\
Geofence radius             & 500 m\\
Incident interval (DES)     & 20--30 s\\
Verification time (DES)     & 1--2 s\\
Dispatch time (DES)         & 1--5 s\\
Failure probability         & 2\,\% (nominal), 20\,\% (failure mode)\\
Independent runs per scen.  & 30\\
Random seed                 & 42 (fixed for reproducibility)\\
\bottomrule
\end{tabular}
\end{table}

% ─────────────────────────────────────────────────────────────
\section{Metrics Evaluated}

The following metrics were collected and analyzed during execution:

\begin{itemize}
  \item Number of processed incidents
  \item Number of verified incidents
  \item Number of failed alert deliveries
  \item Average response time (incident creation to dispatch)
  \item Alert broadcast count (number of users reached)
  \item Latency compliance percentage (incidents resolved in $< 30$~s)
\end{itemize}

For the behavioral simulation, additional metrics were tracked:
community adoption rate, alert fatigue level, user trust level, and
adoption dynamics across the 12-step window.

% ─────────────────────────────────────────────────────────────
\section{Scenario Design and Experimental Matrix}

Three scenarios were defined to test the system under different
operating conditions, aligned with the risk scenarios identified in
Workshop~3.

\begin{table}[H]
\caption{Factorial Experiment Design (EXP-01, EXP-02, EXP-03)}
\label{tab:exp}\centering\small
\begin{tabular}{p{0.6cm}p{2.5cm}p{1.7cm}p{2.2cm}}
\toprule
\textbf{EXP} & \textbf{Risk Hypothesis} &
\textbf{Independent Variable} & \textbf{Dependent Variable (Target)}\\
\midrule
EXP-01 &
  Massive traffic peaks degrade notification latency, causing delayed
  alerts. &
  Incoming report load (100--300\,\% above baseline). &
  Dispatcher latency\\(target: $< 30$~s).\\[4pt]
EXP-02 &
  Report validation saturates security staff, increasing response
  time beyond the quality threshold. &
  Verification complexity during peak hours (18:00--04:00). &
  End-to-end verification latency (target: $< 60$~s).\\[4pt]
EXP-03 &
  Low reporting culture prevents reaching the critical user mass
  needed for reliable detection. &
  UI friction level in the mobile reporting flow. &
  Community adoption rate (target: $\geq 60\,\%$); report
  completeness (target: $\geq 80\,\%$).\\
\bottomrule
\end{tabular}
\end{table}

Each experiment follows the same protocol: fixed random seed
initialization, 12 simulation steps, per-step metric collection,
and paired baseline vs.\ optimization runs to measure the improvement
produced by each mitigation mechanism.

% ─────────────────────────────────────────────────────────────
\section{Extended Risk Register}

Table~\ref{tab:risks} extends the Workshop~3 risk register by
adding the experiment that validates each risk (EXP column) and the
specific technological mechanism from Workshop~2 that acts as the
mitigation.

\begin{table*}[t]
\caption{Extended Risk Register with Experimental Validation (R1--R10)}
\label{tab:risks}\centering\small
\begin{tabular}{clcclllp{3.5cm}}
\toprule
\textbf{ID} & \textbf{Risk} & \textbf{Prob.} & \textbf{Impact} &
\textbf{Level} & \textbf{EXP} & \textbf{WS2 Mechanism} & \textbf{Quality Attr.}\\
\midrule
R1  & System overload        & High   & High   & Critical & EXP-01 &
  Kubernetes auto-scaling + RabbitMQ priority queues & Scalability\\
R2  & Notification failure   & Medium & High   & High     & EXP-01 &
  SMS + Push redundancy; async retry queues & Reliability\\
R3  & False reports          & High   & Medium & High     & EXP-02 &
  Geolocation verification engine; threshold $\geq 3$ & Integrity\\
R4  & Data breach            & Low    & Critical & High   & ---    &
  TLS 1.3 end-to-end; AES-256 at rest & Security\\
R5  & Low adoption           & Medium & High   & High     & EXP-03 &
  Mobile UX $\leq 3$ taps (NFR-06); progressive onboarding & Usability\\
R6  & Connectivity loss      & High   & Medium & High     & EXP-01 &
  Automatic SMS fallback; offline sync & Availability\\
R7  & API integration failure & Medium & High  & High     & EXP-02 &
  Redundant API + circuit breaker; exponential retry & Reliability\\
R8  & Slow authority response & Medium & High  & High     & EXP-02 &
  Automated escalation; monitored SLA & Performance\\
R9  & Poor user experience   & Medium & Medium & Medium  & EXP-03 &
  Iterative UX testing; A/B flow testing & Usability\\
R10 & Database failure       & Low    & Critical & High  & ---    &
  Multi-zone replication; backups every 15 min & Reliability\\
\bottomrule
\end{tabular}
\end{table*}

% ─────────────────────────────────────────────────────────────
\section{Mitigation Strategies}

\subsection{M1 --- Verification Threshold (R1, R2, R3)}

The verification threshold $\theta$ is the most critical control
variable in the system. With $\theta = 1$, any report---even a
false one---triggers a broadcast, which degrades community trust by
approximately 15 percentage points per simulation cycle. With
$\theta \geq 3$, the system filters false positives without
exceeding the latency bound:
\[
\text{latency} = \frac{\text{reports} \times 15}{\text{efficiency}}
< 60\text{ s at efficiency} = 0.9
\]
EXP-01 validates that the dispatcher maintains latency below 30~s
at 300\,\% load when RabbitMQ and Kubernetes are active. EXP-02
confirms that the automated verification engine keeps end-to-end
latency below 60~s.

\subsection{M2 --- Critical Adoption Mass and Minimal UI Friction
  (R5, R9)}

The behavioral simulation shows that the system needs a minimum of
30--35\,\% community adoption to function reliably. Below that
threshold, the human sensor network is too sparse to detect incidents
consistently. EXP-03 tests how simplifying the reporting UI (reducing
steps to $\leq 3$) affects adoption. Progressive onboarding and
social-proof mechanisms help overcome the initial resistance to
installing and using the app.

\subsection{M3 --- Alert Fatigue Prevention (R6, R8)}

The Failure Mode scenario reveals that alert fatigue is irreversible
once it crosses 60\,\%: adoption drops 4 percentage points per
step and the system cannot recover on its own. The threshold
$\theta \geq 3$ is the main mechanism to prevent this collapse.
Automatic SMS fallback covers connectivity failures (R6), and an
automated authority escalation protocol with a monitored SLA
handles slow response times (R8).

\subsection{M4 --- Infrastructure Security and Resilience (R4, R7,
  R10)}

Some risks do not depend on simulation behavior---they require
architectural guarantees. TLS 1.3 and AES-256 protect data in
transit and at rest (R4). Circuit breakers with exponential retry
prevent failures from cascading across microservices (R7). Multi-zone
database replication with 15-minute backups ensures data availability
in case of hardware failure (R10).

% ─────────────────────────────────────────────────────────────
\section{Contingency Plan}

Table~\ref{tab:contingency} defines the immediate response for the
five highest-impact failure scenarios identified during the simulation.

\begin{table}[H]
\caption{Contingency Responses for Top Failure Scenarios}
\label{tab:contingency}\centering\small
\begin{tabular}{p{1.4cm}p{1.6cm}p{2.2cm}p{1.5cm}}
\toprule
\textbf{Scenario} & \textbf{Trigger} &
\textbf{Immediate Action} & \textbf{Mechanism}\\
\midrule
Overload (R1) &
  Latency $>30$~s &
  Scale Kubernetes pods; activate priority queues &
  HPA + PagerDuty\\[4pt]
Notif.\ fail (R2) &
  Delivery $< 95\,\%$ &
  Activate SMS fallback channel &
  RabbitMQ DLQ + Twilio\\[4pt]
Alert fatigue (R3) &
  Fatigue $> 60\,\%$ &
  Raise $\theta$ to 4; pause non-critical alerts &
  Real-time feature flag\\[4pt]
Data breach (R4) &
  SIEM anomaly &
  Isolate service; revoke tokens &
  Circuit breaker + secret rotation\\[4pt]
DB failure (R10) &
  Primary unreachable &
  Promote secondary replica &
  Multi-AZ failover $< 30$~s\\
\bottomrule
\end{tabular}
\end{table}

% ─────────────────────────────────────────────────────────────
\section{Results and Analysis}

\subsection{DES Execution Results}

The DES simulation produced variable results across individual runs,
as expected from a stochastic model. Two representative executions
illustrate the range of behavior:

Under moderate load conditions, the system processed 23 incidents, verified 24, had 0
failed alerts, and broadcast alerts to 30 nearby users, with an average response
time of 15.87~s and 100\,\% of incidents resolved within 30~s.

Under higher load conditions, the system processed 29 incidents, had 1 failed alert
delivery, and broadcast to 47 users, with an average response time of 28.14~s and
only 55.17\,\% of incidents within the 30~s threshold.

This variability shows that performance strongly depends on incident
generation frequency and queue congestion---a finding consistent with
the Workshop~1 observations about peak incident hours.

\subsection{Behavioral Simulation --- Statistical Results}

Table~\ref{tab:results} presents the aggregated results from 30
independent runs per scenario. The confidence intervals were
calculated using Student's $t$-distribution with
$t_{29,\,0.025} = 2.045$.

\begin{table}[H]
\caption{Scenario Results --- 30 Independent Runs with 95\,\% CI}
\label{tab:results}\centering\small
\begin{tabular}{lcccc}
\toprule
\textbf{Scenario} & \textbf{Avg Lat.} &
\textbf{CI Low} & \textbf{CI High} & \textbf{Adoption}\\
\midrule
Baseline     & 40.87 s & 38.85 s & 42.89 s & 21.39\,\%\\
Optimization & 17.94 s & 16.92 s & 18.96 s & 67.62\,\%\\
Failure Mode & 71.66 s & 67.72 s & 75.59 s & 31.34\,\%\\
\bottomrule
\end{tabular}
\end{table}

The Optimization scenario achieved the best results in all metrics,
maintaining low latency, high adoption, and stable alert propagation.
The Failure Mode scenario showed severe queue congestion, high alert
fatigue, and reduced user trust---confirming that infrastructure
degradation amplifies system problems non-linearly.

\begin{figure}[H]
\centering
\includegraphics[width=0.97\columnwidth]{latency_by_scenario.png}
\caption{Average latency by scenario with 95\,\% confidence intervals.
  The Optimization scenario achieves a 56\,\% latency reduction
  compared to Baseline, satisfying NFR-01.}
\label{fig:latency}
\end{figure}

\begin{figure}[H]
\centering
\includegraphics[width=0.97\columnwidth]{adoption_evolution.png}
\caption{Community adoption evolution across 12 simulation steps.
  The Optimization scenario maintains a positive reinforcing loop
  that stabilizes participation above the 60\,\% target.}
\label{fig:adoption}
\end{figure}

\begin{figure}[H]
\centering
\includegraphics[width=0.97\columnwidth]{fatigue_behavior.png}
\caption{Emergent alert fatigue behavior. The Failure Mode scenario
  shows exponential fatigue growth that leads to adoption collapse.
  The Optimization scenario keeps fatigue within a safe range.}
\label{fig:fatigue}
\end{figure}

\begin{figure}[H]
\centering
\includegraphics[width=0.97\columnwidth]{scenario_comparison.png}
\caption{Adoption rate vs.\ alert success rate by scenario.}
\label{fig:comparison}
\end{figure}

% ─────────────────────────────────────────────────────────────
\section{Statistical Validation}

Each scenario was executed 30 independent times using controlled
pseudo-random seeds to reduce stochastic bias. Confidence intervals
were computed using the formula:
\[
\bar{x} \pm t_{29,\,0.025} \cdot \frac{s}{\sqrt{30}},
\quad t_{29,0.025} = 2.045
\]

The results show stable metric distributions with acceptable
variability ranges---the narrow confidence intervals for the
Optimization scenario ($\pm 1.02$~s) indicate consistent behavior
across runs.

Sensitivity analysis identified two variables with the greatest
influence on system stability:
\begin{enumerate}
  \item \textbf{Verification threshold}: increasing from 1 to 3
    reduces false-positive propagation by approximately 78\,\%
    and prevents fatigue collapse.
  \item \textbf{Staff efficiency}: increasing from 50\,\% (baseline)
    to 90\,\% (optimization, representing the automated engine)
    cuts average latency roughly in half.
\end{enumerate}

All other parameters had second-order effects that only become
significant at extreme values outside normal operating conditions.

% ─────────────────────────────────────────────────────────────
\section{Complexity and Emergent Behavior}

\subsection{Feedback Loop Structure}

Figure~\ref{fig:feedback} shows the two competing feedback loops
that govern long-term system behavior. The verification threshold
$\theta$ determines which loop dominates.

\begin{figure}[H]
\centering
\includegraphics[width=0.97\columnwidth]{feedback_loops.png}
\caption{System feedback loop structure. Green path (R+): reinforcing
  loop active when $\theta \geq 3$. Red path (B-): balancing loop
  that drives adoption collapse when $\theta = 1$.}
\label{fig:feedback}
\end{figure}

With $\theta \geq 3$, the R+ loop dominates: verified alerts raise
community vigilance, which drives more reports, sustaining the cycle.
With $\theta = 1$, the B- loop dominates: false alerts reach users,
fatigue grows, adoption collapses, and the system cannot recover
without intervention---a behavior consistent with sensitive dependence
on a control parameter in nonlinear dynamical systems~\cite{sterman2000}.

\subsection{Sensitivity to Initial Conditions}

Figure~\ref{fig:sensitivity} shows how the three main metrics respond
to changes in the verification threshold $\theta$.

\begin{figure}[H]
\centering
\includegraphics[width=0.97\columnwidth]{sensitivity_analysis.png}
\caption{Sensitivity analysis: effect of verification threshold
  $\theta$ on average latency, community adoption, and alert success
  rate. The dashed line marks the recommended value $\theta = 3$.}
\label{fig:sensitivity}
\end{figure}

The curves show a sharp transition between $\theta = 1$ and
$\theta = 3$, after which improvements become marginal. This
inflection is the system's bifurcation point: below it the
balancing loop dominates and the system is unstable; at or above
it the reinforcing loop dominates and the system is stable.
Increasing $\theta$ beyond 3 yields diminishing returns while
reducing the number of legitimate alerts broadcast.

\subsection{Computational Complexity}

The dispatcher operates close to $O(n)$ under moderate load.
During overload, priority-queue scheduling introduces behavior
closer to $O(n \log n)$. The three main bottlenecks identified
were: (1)~incident verification latency, (2)~alert dispatch queue
congestion, and (3)~community validation throughput. Kubernetes
auto-scaling and RabbitMQ priority queues mitigated these in the
Optimization scenario, confirming the Workshop~2 design.

% ─────────────────────────────────────────────────────────────
\section{Validation Checklist}

\begin{table}[H]
\caption{Validation Methods and Acceptance Criteria}
\label{tab:validation}\centering\small
\begin{tabular}{p{3.5cm}p{3.9cm}}
\toprule
\textbf{Validation Method} & \textbf{Criterion / Target}\\
\midrule
Behavioral simulation (3 scenarios) &
  Adoption $\geq 60\,\%$, fatigue $< 60\,\%$, trust $> 80\,\%$\\
Load test (EXP-01) &
  Latency $< 30$~s at 300\,\% traffic\\
Integration test (EXP-02) &
  End-to-end verification $< 60$~s\\
UAT / friction test (EXP-03) &
  Adoption $\geq 60\,\%$; completeness $\geq 80\,\%$\\
Security test &
  No critical vulnerabilities (OWASP Top~10)\\
Bifurcation analysis &
  System stable at threshold $\theta \geq 3$\\
\bottomrule
\end{tabular}
\end{table}

% ─────────────────────────────────────────────────────────────
\section{Assumptions and Limitations}

\subsection{Assumptions}

The simulation was built on the following simplifying assumptions,
which represent the gap between the model and a real deployment:

\begin{itemize}
  \item GPS coordinates are randomly generated within a fixed area.
  \item AI verification confidence is sampled from a uniform
    distribution---a real trained model would behave differently.
  \item User behavior follows simplified probabilistic rules.
  \item Network failures occur as independent random events with
    fixed probability.
  \item Users are uniformly distributed across campus, which does not
    reflect the spatiotemporal concentration identified in Workshop~1.
\end{itemize}

\subsection{Limitations}

\begin{itemize}
  \item No real mobile application was implemented or tested.
  \item The AI verification model is probabilistic, not trained on
    real data.
  \item GPS functionality is fully simulated.
  \item Human behavior is approximated statistically, not observed.
  \item Database persistence is temporary (in-memory only).
\end{itemize}

% ─────────────────────────────────────────────────────────────
\section{Source Code Documentation}

The simulation source code is organized following the same modular
architecture as the system design (NFR-07 Maintainability). Two
main scripts implement the two simulation paradigms:

\textbf{W4\_DES.py} implements the DES using SimPy. Each service
is an independent class with a single responsibility. All
configuration constants are declared at the top of the file,
making it easy to change parameters without modifying the logic.

\textbf{W4\_behavioral.py} implements the behavioral simulation using
a step-by-step state update model. The \texttt{CSASSimulation} class
tracks adoption, fatigue, trust, and vigilance across the 12-step
window. Environmental multipliers are applied through the
\texttt{environmental\_multipliers} function based on the hour,
weather, and campus zone.

\begin{lstlisting}[caption={Verification gate and feedback loops
--- W4\_behavioral.py}]
def run_step(self):
    # Verification gate
    if reports >= self.threshold:
        latency = (reports * 15) / self.efficiency
        if latency < 60:
            alert_sent = 1

    # Reinforcing loop (R+): alert verified, no fatigue
    if alert_sent and self.threshold >= 3:
        self.vigilance = self._clamp(self.vigilance + 0.12)
        self.fatigue   = self._clamp(self.fatigue   - 0.05)

    # Balancing loop (B-): fatigue above 60%
    if self.fatigue > 0.60:
        self.adoption = self._clamp(self.adoption - 0.04, 0.05)
\end{lstlisting}

Each class and function includes a docstring describing its
responsibility, inputs, outputs, and behavior. This structure
supports future modifications without affecting other modules.

% ─────────────────────────────────────────────────────────────
\section{Reproducibility and Output Files}

All simulation outputs are exported to structured files to allow
independent verification of the results.

\begin{table}[H]
\caption{Simulation Output Files}
\label{tab:outputs}\centering\small
\begin{tabular}{ll}
\toprule
\textbf{File} & \textbf{Purpose}\\
\midrule
\texttt{baseline\_log.csv}         & Step-by-step baseline execution trace\\
\texttt{confidence\_intervals.csv} & Per-scenario statistics and CI\\
\texttt{sensitivity\_analysis.csv} & Results for threshold sweep ($\theta=1\ldots4$)\\
\texttt{figures/*.png}             & All four performance charts\\
\bottomrule
\end{tabular}
\end{table}

To reproduce the full simulation:
\begin{lstlisting}[language=bash, numbers=none]
cd simulation/
pip install -r requirements.txt
python W4_DES.py          # Process-oriented DES
python W4_behavioral.py   # Behavior-oriented simulation
\end{lstlisting}

% ─────────────────────────────────────────────────────────────
\section{Ethical Considerations}

All data used in the simulation is synthetic. No personal information
from real students or staff was collected or processed. Any future
production deployment of the CSAS must comply with Ley~1581 de~2012,
the Colombian personal data protection law, particularly regarding
consent for location data collection.

The verification algorithm does not use demographic variables, but
a real trained AI model should be audited for potential bias before
deployment. From a safety perspective, false-negative alerts (a real
incident not broadcast) carry a higher social cost than false
positives, so the threshold $\theta \geq 3$ was chosen to balance
filtering quality against the risk of suppressing legitimate reports.

% ─────────────────────────────────────────────────────────────
\section{Repository Management}

All workshop artefacts are organized in the
\texttt{Workshop\_4\_Simulation/} folder within the course repository
(\url{https://github.com/filipghs/Systems-Analysis}):

\begin{small}
\begin{verbatim}
Workshop_4_Simulation/
  docs/
    Workshop4_CSAS.pdf
    Workshop4_CSAS.tex
  simulation/
    W4_DES.py
    W4_behavioral.py
    requirements.txt
  results/
    figures/  (4 PNG charts)
    confidence_intervals.csv
    sensitivity_analysis.csv
    baseline_log.csv
  risk_management/
    risk_analysis_W4.md
  validation/
    statistical_validation.md
  README.md
\end{verbatim}
\end{small}

% ─────────────────────────────────────────────────────────────
\section{Design Validation}

This section explicitly assesses how the simulation results validate
or challenge the design decisions made across Workshops 1, 2, and 3.

\textbf{Workshop~1 --- Empirical Findings:}
Workshop~1 measured a 35\,\% incident under-reporting rate and
identified 12\,\% as the initial adoption level, with parking lots
and building entrances during 18:00--04:00 as the highest-risk zones.
The simulation confirms that 12\,\% adoption is insufficient for the
system to sustain itself: the behavioral model shows that the
reinforcing loop cannot overcome noise below 30--35\,\% participation.
The time-of-day multipliers in the behavioral simulation (factor
$2.5\times$ at night) directly reflect those Workshop~1 observations.

\textbf{Workshop~2 --- Architecture Decisions:}

\begin{itemize}
  \item \textbf{RabbitMQ + Kubernetes:} \checkmark Validated by EXP-01.
    The Optimization scenario maintains latency below 30~s at 300\,\%
    load, confirming that the event bus and elastic scaling are adequate.
  \item \textbf{Automated Verification Engine:} \checkmark Validated by
    EXP-02. Staff efficiency of 90\,\% (automated) keeps end-to-end
    latency within NFR-01 bounds; 50\,\% (manual) does not.
  \item \textbf{SMS/USSD fallback:} \checkmark Validated by the failure
    probability sweep. The fallback prevents total outage when the
    primary data channel fails.
  \item \textbf{$\leq$3-interaction mobile flow:} \checkmark Validated
    by EXP-03. Reducing UI friction is the primary mechanism for
    crossing the 30--35\,\% adoption threshold.
  \item \textbf{Initial 12\,\% adoption assumption:} $\times$ Challenged.
    The system cannot sustain the reinforcing loop starting from 12\,\%.
    A structured onboarding campaign is required before Phase~1 exit.
\end{itemize}

\textbf{Workshop~3 --- Risk and Quality Framework:}
The simulation provided quantitative validation for 8 of the 10
registered risks. R1 (Critical) and R6 (High) are the most directly
observable: the Failure Mode scenario reproduces their effects
(latency 71.66~s, adoption collapse) exactly as the risk register
predicted. The two structural risks---R4 (data breach) and R10
(database failure)---could not be modeled stochastically but remain
as architectural invariants that bound the validity of the results.

% ─────────────────────────────────────────────────────────────
\section{Recommendations}

Based on the simulation results, the following improvements are
recommended for the next stages of the CSAS project.

\subsection{Short-term (before Phase~1 deployment)}

\begin{enumerate}
  \item \textbf{Set verification threshold to $\theta = 3$} as the
    production default and expose it as a runtime feature flag so
    the operations team can adjust it without redeploying the service.
  \item \textbf{Run a targeted onboarding campaign} to push community
    adoption past 35\,\% before going live. Below that threshold the
    system cannot self-sustain, regardless of how well the technology
    is configured.
  \item \textbf{Implement a fatigue monitoring dashboard} that tracks
    alert fatigue in real time and automatically raises $\theta$ when
    fatigue approaches 50\,\%, preventing the irreversible collapse
    observed in the Failure Mode scenario.
\end{enumerate}

\subsection{Medium-term (Phase~2 improvements)}

\begin{enumerate}
  \item \textbf{Replace the probabilistic AI gate} with a trained
    binary classifier using real incident data. The current model
    assumes a uniform confidence distribution, which overestimates
    performance for edge cases.
  \item \textbf{Use real campus GPS data} to replace uniform random
    coordinate generation. Workshop~1 identified specific high-risk
    zones; incorporating that spatial concentration would produce
    more accurate capacity estimates.
  \item \textbf{Extend the simulation window} beyond 12 steps to
    model seasonal variation (exam periods, holidays) and evaluate
    whether the 60\,\% adoption KPI is sustainable over months.
\end{enumerate}

\subsection{Long-term (future research)}

\begin{enumerate}
  \item Explore agent-based modeling as a complement to the DES,
    assigning individual behavior profiles to each user rather than
    treating adoption as a single aggregate variable.
  \item Investigate integration with the SIURE platform at
    Universidad Distrital to evaluate whether bidirectional data
    exchange creates additional feedback loops not captured here.
  \item Study the effect of gamification mechanics (reputation
    scores, community badges) on the reinforcing loop coefficient
    as a way to accelerate organic adoption growth.
\end{enumerate}

% ─────────────────────────────────────────────────────────────
\section{Conclusions}

The CSAS simulation confirmed that the proposed system design is
feasible under realistic operating conditions. Three conclusions stand
out:

First, the \textbf{verification threshold is the most important
parameter} in the entire system. Setting $\theta \geq 3$ keeps the
reinforcing loop active, adoption above 60\,\%, and latency well
within NFR-01 bounds. Setting $\theta = 1$ drives the system into
a collapse pattern that cannot recover without external intervention.

Second, \textbf{technology alone is not sufficient}. The simulation
showed that below 30--35\,\% community adoption, even a perfectly
configured system cannot sustain reliable operation because the human
sensor network is too sparse. A structured onboarding campaign is
a prerequisite for Phase~1 deployment, not an optional enhancement.

Third, \textbf{the Optimization scenario validates the Workshop~2
architecture}. RabbitMQ priority queues, Kubernetes auto-scaling,
the automated verification engine, and the SMS fallback channel each
played a measurable role in achieving the NFR-01, NFR-02, and
NFR-05 targets under the simulated conditions.

Future work should address the current limitations: replacing the
probabilistic AI gate with a real trained model, using actual campus
GPS data for a more realistic incident distribution, and extending
the simulation to cover the full phased deployment defined in
Workshop~3.

% ─────────────────────────────────────────────────────────────
\begin{thebibliography}{9}

\bibitem{simpy}
Team SimPy, \textit{SimPy: Discrete-Event Simulation for Python},
version~4.1, 2023. Available: \url{https://simpy.readthedocs.io}

\bibitem{law2015simulation}
A.~M. Law, \textit{Simulation Modeling and Analysis}, 5th~ed.
New York: McGraw-Hill, 2015.

\bibitem{sterman2000}
J.~D. Sterman, \textit{Business Dynamics: Systems Thinking and
Modeling for a Complex World}. Irwin/McGraw-Hill, 2000.

\bibitem{iso25010}
ISO/IEC, \textit{Systems and Software Quality Models},
ISO/IEC~25010:2011.

\bibitem{pmbok7}
Project Management Institute, \textit{PMBOK Guide}, 7th~ed.
Newtown Square, PA: PMI, 2021.

\bibitem{iso31000}
ISO, \textit{Risk Management --- Guidelines}, ISO~31000:2018.

\bibitem{dennis2021}
A.~Dennis, B.~H. Wixom, and D.~Tegarden,
\textit{Systems Analysis and Design}, 6th~ed. Hoboken, NJ: Wiley, 2021.

\bibitem{ley1581}
Congreso de Colombia, \textit{Ley~1581 de~2012 ---
Protecci\'on de Datos Personales}, Diario Oficial No.~48.587, 2012.

\end{thebibliography}

\end{document}
